\chapter{Introduction}\label{chap:introduction}

\section{Motivation}\label{sec:motivation}

News moves markets. This is not a new claim. \citet{tetlock2007} showed that pessimism in the financial press predicts downward pressure on prices, and that unusually high or low media sentiment predicts high trading volume. The problem is that most of this work is done on daily data. A daily bar treats the whole session as one number, and the interesting part happens inside it. When a scheduled release lands, the market reacts within minutes to hours and then the reaction fades over the rest of the day \citep{ederington1993}. Daily data cannot tell us whether liquidity moves on impact, an hour later, or only the next day. It has already averaged that away.

West Texas Intermediate (WTI) crude oil is a good place to look at this more closely. The news that drives it is frequent and, more importantly, well dated: Organization of the Petroleum Exporting Countries (OPEC) policy decisions, geopolitics in the producing regions, the weekly U.S. Energy Information Administration (EIA) inventory release, and the usual macroeconomic announcements. Each of these arrives at a time we know, and each could plausibly move liquidity over a few hours rather than a few days. WTI is also deeply liquid, so there is enough trading activity to measure at a fine grain. That combination is what makes an hourly study possible here.

\section{Problem statement}\label{sec:problem-statement}

The question this thesis works on is how news propagates into WTI liquidity at hourly resolution. Three things need answering.

The first is timing. When news arrives, does trading react in the same hour, or later? If later, how much later?

The second is direction. Does bearish news move liquidity more than bullish news of the same size? Loss aversion says it should, so it is worth testing.

The third is context. Liquidity does not respond to news in a vacuum. Some of the response belongs to the news, and some belongs to the state of the market at the time: how volatile things already were, where the dollar was, what kind of event it was. We want to know how that splits, and we want to read the split off a model rather than assume it.

Getting to any of this means two things. Unstructured news text has to become numbers on an hourly grid lined up with market data. Then the news-liquidity relationship has to be modelled with methods transparent enough that the timing and direction results can actually be interpreted, rather than just reported.

\section{Research questions}\label{sec:research-questions}

\begin{description}
  \item[RQ1 (lag structure).] At what lag, if any, does news sentiment exert its strongest effect on WTI trading liquidity, measured at hourly resolution?
  \item[RQ2 (directional asymmetry).] Does the liquidity response differ between bearish and bullish news, and if so, in what sense?
\end{description}

\section{Contributions}\label{sec:intro-contributions}

There are three contributions defined in this thesis. Chapters~\ref{chap:experiments} and~\ref{chap:discussion} develop them deeper.

\textbf{(a) The intraday shape of the news-liquidity response, and its directional asymmetry.} Most news-and-markets work runs at daily resolution. That is enough to establish \emph{that} news moves trading activity \citep{tetlock2007}, but not \emph{when} inside the session, because a daily bar has averaged the answer away before the analysis begins. The intraday work that does exist looks mainly at scheduled macroeconomic announcements, and at prices and volatility, where most of the adjustment completes within minutes of release \citep{ederington1993}. This thesis measures a different object: how \emph{liquidity} responds to a continuous stream of largely unscheduled news, hour by hour.

Working at that resolution is what makes three results visible, none of which a daily design can represent:

\begin{itemize}
  \item \textbf{The response is not contemporaneous.} It peaks at $+6$ hours and has decayed to insignificance by $+8$, so the entire lifetime of the effect fits inside a single daily bar (Section~\ref{sec:lag-ols}).
  \item \textbf{The decay is not smooth.} Secondary maxima at $+1$h and $+4$h point to discrete information-incorporation events rather than a continuous fade.
  \item \textbf{The timing itself is direction-dependent.} Bearish news is integrated over roughly six hours against one hour for bullish, an asymmetry that lives in \emph{when} the market responds rather than in how much (Section~\ref{sec:tft-v2-lag}).
\end{itemize}

Two methods with no assumptions in common agree on the $+6$h peak: a lag-regression baseline (Section~\ref{sec:lag-ols}) and the attention pattern of a Temporal Fusion Transformer (Section~\ref{sec:tft-v2}). Read against \citet{ederington1993}, where prices finish adjusting within minutes, the substantive claim is that \textbf{liquidity is the slower variable}: the market has repriced long before it has finished trading on the news. Reaching that result also meant solving an alignment problem specific to sub-daily work, since news publishes continuously while trading does not, so every article has to be placed on a grid with overnight and weekend gaps in it (Section~\ref{sec:temporal-alignment}).

Asymmetry is more involved. The two phases do not measure the same quantity, so their results look like a disagreement at first and are not one. The linear baseline gives a larger marginal coefficient on negative sentiment than on positive (Section~\ref{sec:lag-ols}). The forecasting model finds no bearish-versus-bullish gap in predicted volume at all, and instead leans on what the news implies about risk: market volatility and the extracted supply channel are its two highest-weighted features, while the composite sentiment score ranks 21st of 90 (Section~\ref{sec:tft-v2}). Put together, the conclusion is not that bearish news beats bullish news. It is that liquidity cares more about how much risk the news carries than about which way it points.

\textbf{(b) A channel decomposition of LLM-extracted news features, checked by cross-model agreement.} Collapsing an article to one sentiment score throws information away. Phase~2 instead extracts separate supply, demand, and risk-premium channels with a large language model, following the structural decomposition of oil-price shocks in \citet{kilian2009}. There is no human ground truth to check the extraction against, so we check it a different way: two models are run over the same articles, and we ask how much they agree. Recalibrating onto the channels lifts cross-model correlation on the composite sentiment from 0.39 to 0.88 (Section~\ref{sec:calibration}). The decomposition also clears up something that looks odd otherwise, which is why supply-risk geopolitical news comes out price-bullish rather than bearish.

\textbf{(c) A comparison of title-only and full-body news representations.} Sentiment taken from headlines alone disagrees with sentiment from full article bodies on 41.6\% of articles. The disagreement is not random either: when an article flips, it flips bearish far more often than bullish (Section~\ref{sec:headline-bias}). Studies that use titles alone are picking up this bias whether they intend to or not, and this puts a number on it.

There is a common thread. Phase~1 is a simple, interpretable baseline (FinBERT features, lag regression). Phase~2 is richer and more expressive (LLM channels, a Temporal Fusion Transformer). The second corroborates the first and gives us more to say about it, and neither one gets thrown away (Section~\ref{sec:methodological-contributions}).

\section{Structure of the thesis}\label{sec:thesis-structure}

On Chapter~\ref{chap:background}: Backgroud, we remind the terminology and places this thesis uses against prior work on news and markets, financial-text models, and interpretable forecasting. On Chapter~\ref{chap:methods}: Methods, we then cover the two-phase methodology: getting the data and aligning it to an hourly grid, extracting features by phases (FinBERT in Phase~1, LLM channel extraction in Phase~2), and the modelling and evaluation setup. On Chapter~\ref{chap:experiments}: Experiments we report the experiment, what the headline-bias comparison means and the lag regression from Phase~1, following the Temporal Fusion Transformer results from Phase~2 with its feature importance, temporal attention, and the directional-asymmetry test. Later on Chapter~\ref{chap:discussion}: Discussion, reads the findings back against the research questions and sets out the limitations found in the experimentation phase. Finally, Chapter~\ref{chap:conclusion}: Conclussion closes the thesis giving out what we learnt and what we could work on in the future. After the main chapters, the appendices hold the extraction prompt and schema, the ablation and hyperparameters, the canonical entity list, and a reproducibility statement mentioned through the previous chapters.
